\section*{September 22, 2026: Flagging REO resales and price tails instead of removing them}
\textit{Drafted by Claude; Jacob decided to flag REO resales and keep a single tail rule.}

The cleaner now removes only sales whose recorded price is not a market price.
These are foreclosure auctions and transfers to a lender or land bank: about
200 to 830 a year in 2008--2018. Resales by lenders, servicers, Fannie Mae,
Freddie Mac, HUD or the land bank are real sales of distressed homes, and
foreclosure incidence may itself respond to closures. They are therefore kept
and flagged \texttt{reo\_sale}. REO resales are 35\% of 2008--2012 sales and
18\% afterwards. Dropping them shapes the pre-period sample, so estimates
should be shown with and without them. The earlier entry's name rules still
define these sales; only the treatment of REO resales changed.

The within-year 99.9th-percentile price-per-square-foot trim was removed, which
restores 133 sales. The hard error screens (fewer rooms than bedrooms, over
\$5,000 per square foot) remain. The single tail rule is now the flag for the
within-year citywide 1st--99th price percentiles, computed on the sample that
includes REO resales. A new flag, \texttt{same\_party\_names}, marks 107 sales
whose normalized buyer and seller names are identical, excluding land trustees.
This matches the related-party rule in the aldermanic-privilege project. Surname
matching was rejected as unreliable. The citywide 2008--2018 sample is 167,312 sales. It
equals the previous 126,498 plus the 40,681 REO resales and the 133 formerly
trimmed sales. No regression has been rerun.

\small
Reproduction: run \texttt{make} in \texttt{tasks/clean\_home\_sales/code}; the
build prints annual removal and flag counts.
\normalsize
